% Source: https://artofproblemsolving.com/community/c3h2519310_really_hard_thingy
% Stable identity: aops-topic-2519310-q01
\textbf{(a)} For $x>0$, $x^x=e^{x\ln x}$, so it suffices to show $x\ln x\to0$. Put $x=e^{-t}$ with $t\to+\infty$; then $x\ln x=-te^{-t}$. For $t>0$, $e^t\ge1+t+\dfrac{t^2}{2}>\dfrac{t^2}{2}$, so
\[0<te^{-t}<t\cdot\dfrac{2}{t^2}=\dfrac2t\to0.\]
Hence $x\ln x\to0$ and, since $e^u$ is continuous at $u=0$, $x^x\to e^0=1$. (Numerically, $\left(\dfrac{1}{1000}\right)^{\dfrac{1}{1000}}=e^{-\dfrac{\ln1000}{1000}}\approx e^{-0.0069}\approx0.993$.)

\textbf{(b)} As $x\to0^+$, $-\dfrac1x\to-\infty$, so $f(x)=e^{-\dfrac1x}\to0$. But
\[f(x)^x=\left(e^{-\dfrac1x}\right)^x=e^{-1}=\dfrac1e\]
for every $x>0$, so the limit is $\dfrac1e$.

\textbf{(c)} In (a) and (b) the base and the exponent both tend to $0$, yet the powers tend to $1$ in one case and to $\dfrac1e$ in the other. (Replacing $-\dfrac1x$ by $-\dfrac cx$ gives the limit $e^{-c}$, any number in $(0,1)$.) So knowing that ``base $\to0$ and exponent $\to0$'' does not determine the limit of the power: $0^0$ is an \emph{indeterminate form}. The limit in (a) is one particular path and does not single out a value. The convention $0^0=1$ is used in algebra and combinatorics (for example in $(1+x)^n=\sum\binom nkx^k$ at $x=0$, or because there is exactly one function from the empty set to the empty set), but it is a definition chosen for convenience, not a consequence of the limit of $x^x$.

$\boxed{x^x\to1,\ \text{but}\ \left(e^{-\dfrac1x}\right)^x\to\dfrac1e;\ 0^0\text{ is not determined by limits}}$
